\appendix
\section{Realization and archival crosswalk for every finding}

The following anchors locate the mathematical framework of a
finding, not its theta formula in Grothendieck's handwriting.
Write G44 for cote~44, scans~7--8 (12 August 1964), and G16 for
cote~16, scans~43--46 (4 January 1967). The original scans are
\url{https://grothendieck.umontpellier.fr/44.pdf} and
\url{https://grothendieck.umontpellier.fr/16.pdf}.
The published \emph{Standard conjectures on algebraic cycles},
pp.~196--198, is an additional polarization/realization anchor.

\begin{longtable}{@{}p{.18\linewidth}p{.48\linewidth}p{.27\linewidth}@{}}
\toprule Finding IDs & Geometric or motivic meaning & Exact framework anchor\\\midrule\endhead
A1--A3 & Oscillator period functionals on a Betti realization; extraction before an algebraic image cycle. & G44, scans 7--8.\\
A4 & Restriction to ample complete intersections; compatibility of the cup pairing. & G16, scans 43--45.\\
A5 & Positive analytic recovery operator on canonical forms; algebraicity is a separate question. & G44, scans 7--8; G16, scans 44--45.\\
A6 & Actual divisor contractions and image-cycle pushforwards, hence algebraic correspondences modulo homological equivalence. & G44, scans 7--8; G16, scans 43--46.\\
A7 & Exact analytic theta realization; its source vector is not asserted to be a motivic correspondence. & G44, scans 7--8.\\
B0--B3 & Hermitian realization data and the local representation needed for a period comparison. & G44, scans 7--8, as the eventual correspondence setting.\\
B4--B8 & Positivity, determinant and overlap estimates after fixing realization norms. & G44, scans 7--8; published Standard conjectures, pp.196--198.\\
B9--B12 & Local and global definite-subgroup coefficients of the canonical line; evaluation is not pullback to degree-two cohomology of a point. & G44, scans 7--8.\\
B13 & Uniform functional on an iterated theta family, with arithmetic links retained. & G44, scans 7--8; G16, scans 44--45.\\
B15 & Explicit local theta representations, global splitting characters and exact normalized two-step kernels. & G44, scans 7--8; G16, scans 44--45.\\
B14 & Algebraic diagonal contraction of specified tensor realizations, contrasted with a different coherent line-bundle power. & G16, scans 43--46.\\
Q1--Q3 & Quantitative selection of forms and rational multiplicity data preceding cycle construction. & G44, scans 7--8.\\
Q4--Q8 & Explicit normalized period and polarization certificates; Q4,Q5 are aliases of A2,A3. & G44, scans 7--8; published Standard conjectures, pp.196--198.\\
Q9 & Algebraic endomorphism polynomials; idempotence modulo homological equivalence, with denominators. & G16, scans 43--45.\\
Q10 & All-degree determinant and period certificates for the A6 cycle construction. & G44, scans 7--8; G16, scans 43--46.\\
C01--C03 & Albanese realization, trace projector and genuine non-CM factors on noncongruence covers. & G44, scans 7--8.\\
C04--C05 & Correspondence-module extraction and its Hodge-group scope. & G44, scans 7--8; G16, scans 43--45.\\
C06--C08 & Derived local intersections and the horizontal/special-fiber distinction. & SGA4$\frac12$, Cycle, Rem.2.3.9 p.150; C17's exact specialization locators.\\
C09--C11 & Local deformation and derivative valuations; an integral frame is additional input to a cycle comparison. & EGA IV$_2$, Thm.6.9.1 p.153; G44 scans 7--8 as the eventual realization setting.\\
C12 & Image cycle with degree500; product of diagonal divisors. & G44, scans 7--8; G16, scans 43--45.\\
C13--C14 & Explicit Weil projector and its action on the Lefschetz algebra. & G16, scans 43--46; G44, scans 7--8.\\
C15--C16 & Albanese quotients from finite \'etale covers and exact image-cycle degrees. & G44, scans 7--8; G16, scans 43--46.\\
C17 & Relative algebraic correspondences and specialization. & EGA IV$_3$, 8.8.2--8.8.3 pp.28--34; SGA4 III, XVI 2.1 p.213, 2.3 p.215; SGA4$\frac12$, Cycle 2.3.2--2.3.9 pp.145--150.\\
C18 & Integral line-bundle ratios on an explicit tame isogeny cover; connection data retained. & EGA IV$_2$, 6.9.1 p.153; SGA4$\frac12$, Cycle 2.3.2 p.145, 2.3.8 p.148.\\
\bottomrule
\end{longtable}

\section{Source register and exact version conventions}

All cited supplied works remain in the package under their original
filenames. In particular, 032 means
\emph{The rational Hodge conjecture for CM abelian varieties}; its
printed theorem/page locators are used throughout. The companion
van Hoften--Sempliner paper is the supplied Hexagon 2610.00259v2.
Milne's supplied Hexagon 2611.00244v1, Theorems~2.1, 3.1 and~3.5,
is carried as the stated CM-Hodge, finite-field Tate and Hodge-standard
background. It is distinct from his 1999 article and modular-forms notes.

The supplied Markman source is \emph{Cycles on abelian $2n$-folds of
Weil type from secant sheaves on abelian $n$-folds},
\path{LIT_Markman_Hodge_abelian_fourfolds_2502.03415/250531.tex}.
Its theorem labelled \texttt{thm-algebraicity} in Section~1.5 and
corollary in Section~1.6 cover the stated Weil-sixfold and
abelian-fourfold results. These are background when determining
whether a proposed non-CM target is already known; this run does
not claim to improve those general results.

The external primary-source URLs and exact reading locators are
listed in \path{documentation/external_sources_C.tsv} and the
dated source-check records. Important additional versions are:
\begin{itemize}[leftmargin=*]
\item BMM: Bergeron--Millson--M\oe glin, published 2016 version,
125 PDF pages; Lemmas~5.5, 5.18, Propositions~5.17, 5.19,
Theorem~9.1. The 97-page preprint has different page numbers.
\item Liu: \emph{Fourier--Jacobi cycles and arithmetic relative
trace formula}, \url{https://yhzhumath.github.io/FJ.pdf};
Definitions~4.3, 4.5, 4.12; Theorems~4.15, 4.18;
Corollary~4.20; Appendix~D.
\item Haan: arXiv:1403.5061v7, Theorem~3.4, p.8,
\url{https://arxiv.org/pdf/1403.5061v7}.
\item Gan--Qiu--Takeda: arXiv:1207.4709v3,
Sections~11.3--11.7; its stated dimension range is retained.
\item Gross: \emph{Incoherent definite spaces and Shimura varieties},
25 February 2018, Theorem~3.1, pp.6--7.
\item Milne: \emph{Lefschetz classes on abelian varieties},
\url{https://jmilne.org/math/articles/1999aP.pdf},
Proposition~5.4 and Corollary~5.5, p.25, used in C14.
\item Katagiri--Kobayashi: arXiv:2609.07024v1,
Definition~3.1, Remark~3.2 and Lemma~3.3, pp.9--11,
\url{https://arxiv.org/pdf/2609.07024}.
\end{itemize}

Alp\"oge locators use the pinned mirror commit given at the start
of the paper. The full relative file path, declaration, and where
available line numbers are recorded with the relevant finding.
\path{documentation/alpoge_mirror/pinned_manifest.json} and
\path{documentation/step2_archive_manifest.json} preserve recovery
URLs and hashes. A failed retrieval is explicitly recorded rather
than represented as a source reading. No source-module text was
executed, and no Lean build certificate is claimed.

The files \texttt{CHECK\_EVENTS.tsv}, \texttt{QUESTIONS.tsv},
\texttt{LEDGER.md}, \texttt{IMPLICATIONS.md} and
\texttt{TRANSCRIPTION\_NOTES.tsv} are part of the mathematical
deliverable. An empty transcription register means that no
supplied text-of-record reading was available for a justified
Grothendieck correction; it is not a claim that all possible
transcriptions are correct.
